\chapter{The Legacy Engine}
łabel{ch:legacy−engine}
∈dex{legacy}∈dex{succession}∈dex{faction building}

\begin{tcolorbox}[colback=blue!5!white,colframe=blue!75!black,title=Planning for Legacy]
∈dex{legacy!planning}
If you are playing a long campaign or enjoy the idea of your character's influence outlasting them, consider investing in a \textbf{Follower} (Cap 3+) early. That follower can become your \textbf{successor} when your character retires or falls −− inheriting your faction standing, one unresolved burden, and a Legacy Bond. See the checklist at the end of this chapter for steps to prepare.
\end{tcolorbox}

\section*{Succession & Faction Building}
\addcontentsline{toc}{section}{Succession & Faction Building}
∈dex{succession!overview}∈dex{faction building!overview}

\begin{quote}
\emph{``Your legacy is not what you build −−− it is who you build.''}
\end{quote}

The Legacy Engine is a narrative framework that helps your campaign grow across generations without losing emotional continuity. It treats succession as a story event −−− earned through action, consequence, and choice −−− rather than a mechanical transaction.

This chapter is for tables where characters retire, die, or pass on their cause to a follower. It is not required for every campaign, but when the moment comes, these tools help the transition feel earned and the new character feel connected.

\section{Core Philosophy}
łabel{sec:legacy−philosophy}
∈dex{legacy!philosophy}

The Legacy Engine follows three principles:

\begin{enumerate}
  ℹtem \textbf{Narrative Primacy} −−− Succession happens because the story demands it, not because a timer fills.
  ℹtem \textbf{People Before Power} −−− A follower becomes a PC because of who they are, not what they can do.
  ℹtem \textbf{Consequences Persist} −−− Every legacy creates debts, rivals, and unfinished work that the successor must carry.
\end{enumerate}

Legacy is not inherited cleanly. It comes wrapped in obligation.

\section{When Does Succession Happen?}
łabel{sec:succession−triggers}
∈dex{succession!triggers}

Succession is not automatic. It emerges from play. Common triggers include:

\begin{itemize}
  ℹtem A PC retires after completing a major arc.
  ℹtem A PC dies or is permanently removed from play.
  ℹtem A player wants to explore a new character but keep the same faction or cause alive.
  ℹtem The campaign shifts to a new generation (years or decades later).
\end{itemize}

In all cases, the transition should be discussed openly at the table. The GM and player collaborate to find the right moment.

\section{Choosing a Successor}
łabel{sec:choosing−successor}
∈dex{succession!choosing a successor}∈dex{followers!as successors}

A successor should already exist in the fiction. Typically, this is a \textbf{Follower} (see §\ref{ch:assets−followers}) with whom the PC has a meaningful relationship. It could also be a family member, a student, a loyal lieutenant, or even a rival who earns respect.

\paragraph{What Makes a Good Successor?}
\begin{itemize}
  ℹtem They have a name, a face, and a voice.
  ℹtem They have acted independently (off‑screen or on‑screen) at least once.
  ℹtem They want something that the PC's faction or cause can provide.
  ℹtem They have a flaw or burden that will complicate their rise.
\end{itemize}

If no suitable follower exists, the GM and player can create one retroactively −−− but that character arrives with less established history. The player may need to spend a session or two building them up before succession feels earned.

\section{The Succession Arc}
łabel{sec:succession−arc}
∈dex{succession!arc}

A smooth transition usually follows three beats. They can span one session or several, depending on the table's pace.

\subsection{1. The Call}
łabel{subsec:succession−call}
∈dex{succession!the call}

The follower faces a moment where they must act beyond their role. Examples:
\begin{itemize}
  ℹtem The PC is incapacitated, absent, or compromised.
  ℹtem A decision must be made that no one else can make.
  ℹtem A crisis demands leadership rather than obedience.
\end{itemize}
This scene answers: \textit{Why this person, now?}

\subsection{2. The Trial}
łabel{subsec:succession−trial}
∈dex{succession!the trial}

The follower is tested. The trial can be:
\begin{itemize}
  ℹtem A single high‑stakes scene (e.g., a duel, a negotiation, a desperate escape).
  ℹtem A short sequence of escalating challenges (2−−3 scenes).
  ℹtem A moral or political decision with lasting impact.
\end{itemize}
Success is not guaranteed. Failure does not end the story −−− it changes it. A follower who fails may still become a PC, but they carry a scar, a doubt, or a broken promise into their new role.

\subsection{3. The Ascension}
łabel{subsec:succession−ascension}
∈dex{succession!the ascension}

At a dramatically appropriate moment −−− often during downtime or immediately after a crisis −−− the follower formally takes up the mantle. This might involve:
\begin{itemize}
  ℹtem A public ceremony (crowning, oath‑swearing, banner‑raising).
  ℹtem A private passing of a symbol (sword, key, ledger).
  ℹtem A quiet acknowledgment (``You're in charge now.'').
\end{itemize}
The departing PC may be present, or they may be gone. Either way, the scene marks the end of one chapter and the beginning of another.

\section{What the Successor Inherits}
łabel{sec:inheritance}
∈dex{succession!inheritance}

When a follower becomes a PC, they do not simply copy the old character sheet. Instead, they inherit three things:

\subsection{1. Faction Standing}
łabel{subsec:inherit−standing}
∈dex{succession!faction standing}

The successor begins with the same reputation, allies, and enemies that the previous PC cultivated. The GM may adjust specific relationships based on how the succession occurred (e.g., a contested succession might lose one ally but gain a bitter rival).

\subsection{2. One Unresolved Burden}
łabel{subsec:inherit−burden}
∈dex{succession!burden}

The player and GM choose one burden that the successor must carry forward. Examples:
\begin{itemize}
  ℹtem \textbf{Oath of Succession}: ``You must protect the valley as I did.''
  ℹtem \textbf{Unfinished Task}: ``Find my daughter. Tell her I'm sorry.''
  ℹtem \textbf{Rival Claim}: Another faction −− or a jealous lieutenant −− disputes your right to lead.
\end{itemize}
This burden is a narrative hook, not a mechanical penalty. It gives the new character immediate purpose.

\subsection{3. A Legacy Bond}
łabel{subsec:inherit−bond}
∈dex{succession!legacy bond}∈dex{bonds!legacy}

The successor has a bond to their predecessor. Once per session, the player may invoke this bond to:
\begin{itemize}
  ℹtem Recall a memory that provides insight.
  ℹtem Receive instinctive guidance in a familiar situation.
  ℹtem Gain a free \textbf{Boon} (see §\ref{ch:core−mechanics}) when acting to uphold the predecessor's values.
\end{itemize}
This bond fades over time as the successor finds their own identity −− but it never vanishes entirely.

\section{Building a Faction Through Succession}
łabel{sec:faction−growth}
∈dex{faction building!through succession}∈dex{legacy!faction growth}

Each time a character passes the torch, the faction (or cause) grows. It accumulates history, traditions, and institutional memory.

\subsection{Growth Through People}
łabel{subsec:faction−people}
Rather than tracking territory or wealth, track the faction's \textbf{reach} through the people who are now part of its story. Each successor adds:
\begin{itemize}
  ℹtem A new contact or ally they brought with them.
  ℹtem A new problem they inherited.
  ℹtem A new custom or saying that arose from their trials.
\end{itemize}
The GM keeps a simple list: \textit{``The Banner of the Broken Wheel,'' ``the annual vigil at Sootfall Abbey,'' ``the debt to the Mistland ferryman.''} These are the faction's real assets.

\subsection{Faction Timers (Optional)}
łabel{subsec:faction−timers}
∈dex{timers!faction}∈dex{faction building!timers}

For tables that want mechanical weight, the faction can track two simple timers:
\begin{itemize}
  ℹtem \textbf{Growth (6−−8)}: When filled, the faction gains a new narrative permission (e.g., a safehouse in a new city, a reliable spy, a permanent shrine).
  ℹtem \textbf{Health (6)}: When filled, a crisis erupts −− a rival attacks, a leader dies, a secret is exposed. The current PC must deal with it or watch the faction fracture.
\end{itemize}
These timers advance only when the fiction demands −− not automatically.

\section{Example: The Fallen Captain}
łabel{sec:example−fallen−captain}
∈dex{legacy!example}

Miranda has played Michael, captain of the Iron Tide mercenary company, for two years. Michael dies defending a village from a wyrm. The player wants to continue the story.

\paragraph{The Successor}
Michael's second‑in‑command, a grizzled scout named Vesper, has been a follower (Cap 3) for many sessions. Vesper once disobeyed an order to save wounded civilians −− a defining moment.

\paragraph{The Call}
At Michael's funeral pyre, the company looks to Vesper. No one else steps forward.

\paragraph{The Trial}
Vesper must negotiate a new contract with a wary noble who doubts the company's future without Michael. The player rolls, succeeds with cost −− the noble agrees, but demands a hostage.

\paragraph{The Ascension}
Vesper takes up Michael's banner. The player erases Michael's character sheet and creates Vesper as a new PC (starting XP based on campaign tier). Vesper inherits:
\begin{itemize}
  ℹtem The Iron Tide as a \textbf{Mercenary Contract (Standard Asset)} (see §\ref{ch:compendium−assets}).
  ℹtem The burden: \textit{``Find the noble's missing son −− or lose the contract.''}
  ℹtem A legacy bond to Michael: once per session, recall Michael's advice to gain a Boon when leading the company.
\end{itemize}
The campaign continues without missing a beat.

\section{Legacy Planning Checklist}
łabel{sec:legacy−checklist}
∈dex{legacy!checklist}

Before your character retires, dies, or steps aside, work with your GM to prepare your successor. Use this checklist to ensure a smooth transition.

\begin{tcolorbox}[colback=green!5!white,colframe=green!75!black,title=Player's Legacy Planning Checklist,fonttitle=\bfseries]
\begin{enumerate}
  ℹtem \textbf{Name a Successor} −− Choose a Follower (Cap 3+), family member, or loyal ally. Give them a name, a face, and a voice.
  ℹtem \textbf{Clarify the Relationship} −− What does your character mean to them? What do they owe each other?
  ℹtem \textbf{Identify One Unresolved Burden} −− What task, oath, or problem will the successor inherit? (e.g., ``Find my stolen sword,'' ``Protect the village,'' ``Pay the guild debt.'')
  ℹtem \textbf{Pass Down an Asset (Optional)} −− If your character controls a Minor or Standard Asset (see §\ref{ch:compendium−assets}), the GM may allow the successor to retain it with suitable narrative justification (e.g., ``The family safehouse,'' ``The guild seat passes to the heir''). Major or Tier III Assets usually require the successor to earn them through play.
  ℹtem \textbf{Settle Outstanding Obligations} −− Clear or transfer any standing debts to terrestrial patrons (see §\ref{ch:world−powers−obligation}) before stepping down. Unpaid debts may become part of the successor's burden.
  ℹtem \textbf{Hold a Ceremony (On‑Screen)} −− A funeral, a passing of the torch, a swearing of oaths. This scene marks the transition and gives the new PC immediate presence.
  ℹtem \textbf{Create the New Character} −− Build the successor using standard character creation rules, adjusted to the campaign's current tier. The GM may award a small XP bonus (e.g., 4−−8 XP) to represent inherited training or resources.
  ℹtem \textbf{Keep a Legacy Note} −− Write one sentence that the old character leaves behind (a proverb, a scar, a rumour). The GM can use it as a future story hook.
\end{enumerate}
\end{tcolorbox}

\section{Campaign End: The Final Legacy}
łabel{sec:final−legacy}
∈dex{legacy!final}

When the campaign ends, the player may describe one lasting mark the faction leaves on the world. Examples:
\begin{itemize}
  ℹtem A new law named after the first PC.
  ℹtem A shrine where the PC fell, still tended by followers.
  ℹtem A phrase that enters common speech (``Stand like Michael'').
\end{itemize}
The story ends −− but the legacy remains.

\begin{quote}
\emph{A character dies. A faction endures. A story echoes forward.}
\end{quote}